\begin{abstract}
Network intrusion detection in Industrial Internet of Things (IIoT) environments is challenged by scarce labels, heterogeneous device behavior, and high-dimensional traffic. This paper presents ContrastiveXAI-FS, an unsupervised feature-refinement framework that combines self-supervised representation learning, latent-space subset search, and cluster-level explainability. A VICReg encoder is trained with group-aware augmentation that masks semantically related traffic features, a bidirectional search removes features according to their ability to preserve cluster separability in the learned latent space, and a Random Forest proxy with Shapley additive explanations characterizes the resulting unlabeled traffic clusters. The complete pipeline is executed independently within each cross-validation training fold. \rev{On the DataSense CIC IIoT 2025 dataset, the method retains on average 59 of 71 features and achieves macro-F1 scores of 0.936 for binary detection and 0.908 for eight-class classification with Random Forest; an equivalence test confirms performance within $\pm$0.01 of the all-feature configuration, and the method significantly outperforms MCFS and the supervised DataSense-17 subset. Results are consistent across three independent master seeds and extend to N-BaIoT, where the method is equivalent to all features and outperforms MCFS. The discovered clusters are stable across folds and seeds (adjusted Rand index 0.94), and the proxy reproduces them with 0.97 balanced accuracy on unseen test folds, supporting faithful behavioral profiles dominated by TCP-flag, time-to-live, and packet-size statistics. Main results use sample-level folds and therefore measure within-dataset generalization; device- and attack-execution-grouped evaluations are reported to quantify this effect.}
\end{abstract}

\begin{keywords}
Self-supervised learning, contrastive learning, intrusion detection systems, unsupervised feature selection, explainable artificial intelligence.
\end{keywords}
